\chapter{Open Problems and Future Directions}

\section*{What the Bottleneck Framing Has Actually Shown}

Fifteen chapters have now traced a single recurring shape. A generation of machine learning hits a wall: exhaustive search doesn't scale, a single layer can't represent XOR, gradients vanish across depth or across time, features have to be hand-chosen, a sequence has to be processed one step at a time. Somebody finds a way past that specific wall, and the resulting system runs headlong into a different one. Backpropagation solved representation and exposed trainability. Depth, once made trainable, exposed feature engineering as the next constraint. Attention solved sequential processing and exposed a quadratic computational cost. Scale solved predictability and exposed a gap between fluent prediction and trustworthy behavior. Nothing about this pattern suggests it stops at the point this book's history catches up to the present. This closing chapter takes stock of where that same pattern is currently running, honestly, without pretending any of the following problems already has the kind of clean resolution the earlier chapters were able to report.

\section*{Data Is Not Infinite}

Every chapter since Chapter 5 has assumed that more, better data is available if a team is willing to pay for it. That assumption has a shelf life. A 2024 analysis by Pablo Villalobos and colleagues estimated the total stock of public, human-generated text suitable for training language models at somewhere on the order of a few hundred trillion tokens, and projected that, if the recent trend of ever-larger training runs continues, that entire stock would be fully absorbed into training somewhere around the second half of this decade, with meaningful uncertainty on the exact date but not on the basic shape of the trend. This isn't a claim that machine learning progress stops once that point arrives. The same paper, and much of the surrounding research community, points to several directions already under active investigation: synthetic data generated by existing models, transfer of capability from data-rich domains like code into data-poorer ones, and simply extracting more capability out of the same amount of data through better training methods. Whether any or all of these fully substitute for genuinely new, human-generated text is not yet settled. What is reasonably well established is that the specific growth pattern of the last decade, training sets ten times larger every few years, drawn from an ever-expanding pool of freely available text, cannot continue indefinitely on its current terms.

\section*{Knowing What's Actually Happening Inside}

Every architecture in this book, once trained, is a large collection of numbers whose individual values are essentially uninterpretable by direct inspection; Chapter 9's single learned filter, readable at a glance because it had nine weights, is not representative of a model with billions of them. Mechanistic interpretability, the effort to understand what a trained network's internal computations actually represent and how they combine to produce a given output, is an active area of research across multiple labs and universities, not a solved problem with an agreed methodology. Some progress has been made in identifying specific, human-interpretable patterns inside trained models, structures that appear to correspond to concepts, or circuits that appear to implement specific operations. Whether this kind of analysis will scale to fully explaining the behavior of the largest models, or will instead remain a partial, spot-check tool applied to specific behaviors of specific concern, is genuinely unresolved. This connects directly back to Chapter 1's sidebar on the Maxima script that turned out to be silently wrong: a system can be exactly, mechanically specified and still be difficult to verify in practice, and a learned system, whose specification is millions or billions of trained numbers rather than a few dozen lines of code, inherits that difficulty at a scale the earlier chapters never had to confront.

\section*{The Cost of Scale Itself}

Chapter 12's scaling laws describe how loss improves with more compute. They say nothing about what that compute costs, in money, in electricity, or in the physical infrastructure required to supply both. Training runs at the frontier of current model scale require specialized hardware, purpose-built data centers, and electricity consumption significant enough that major technology companies have begun securing dedicated power generation for it. Whether continued scaling along the trajectory this book has traced is economically and physically sustainable indefinitely, or whether it runs into its own hard limit well before any of the other constraints in this chapter bind, is an open question this book won't attempt to forecast, beyond noting that it is a real constraint under active discussion, not a hypothetical one.

\section*{Reliability, Reasoning, and the Unfinished Business of Chapter 14}

Chapter 14 was direct about this already, and it belongs in this chapter's accounting rather than being treated as settled. Models trained through supervised fine-tuning and RLHF are substantially more useful and better behaved than raw pretrained models, and the reward-hacking dynamic that chapter demonstrated directly is genuinely mitigated, not eliminated, by current techniques. Models still, on occasion, produce fluent, confident, entirely fabricated claims, a behavior usually called hallucination, that looks structurally similar to this book's very first sidebar on Chapter 1: a system producing output that is well-formed and plausible while being wrong, without any internal signal distinguishing that case from a correct one. Chapter 15's retrieval and tool use narrow the specific cases where this happens by giving a model somewhere to check its answer against, but they depend on a model reliably recognizing when it should use them, which is itself a trained behavior subject to exactly the same imperfect-proxy problem as everything else in Chapter 14. None of this is a reason for despair; it's a reason to track this as a live, actively worked research problem rather than a resolved chapter of history the way Chapter 4's resolution of the XOR problem was.

\section*{Agents, Compounding Error, and Autonomy}

Chapter 15 closed by noting that an agent, a system that takes multiple sequential actions in pursuit of a goal, inherits and compounds whatever unreliability exists at each individual step. A model that answers a single question incorrectly one time in twenty is a different kind of concern than a model that takes twenty sequential actions, each with a one-in-twenty chance of a consequential error, in service of a single task it was given autonomy to pursue. As these systems are given increasing ability to take real-world actions, browsing, executing code, sending messages, operating other software on a person's behalf, the stakes attached to Chapter 14's unresolved reliability problem scale accordingly. This is an active area of both technical research, on how to make individual steps and error-correction across steps more reliable, and of ongoing discussion about what degree of autonomous action is appropriate to grant a system whose individual judgments remain imperfect in ways this book has demonstrated directly, more than once, rather than merely asserted.

\section*{Returning to Where This Book Began}

This book's prologue opened with a single question, forced on Immanuel Kant by his reading of Emanuel Swedenborg: when does internally coherent reasoning count as knowledge, rather than merely fluent invention? It also made a prediction, before any of the technical history had been told: that every mechanism increasing what a system can express would eventually create a new problem of whether that expression can be trusted. Fifteen chapters of specific, verified evidence now stand behind that prediction. Symbolic rules were exact but couldn't scale. Statistical methods handled uncertainty but needed hand-built features. Deep networks learned their own features but needed a mechanism to stay trainable. Attention solved trainability's remaining cost but introduced a quadratic price. Scale bought predictable improvement in fluency without buying trustworthiness, and post-training narrows that gap without closing it, as this book has shown breaking down directly rather than merely asserted. Retrieval and tool use give a model something external to check itself against, the closest thing in this whole history to Kant's original proposed solution, a boundary supplied from outside the generative process itself, rather than trusted to emerge from within it.

Every chapter's resolution, looked at this way, was really a relocation. That is the honest shape of this book's whole history, not a discouraging one. Each specific bottleneck in each specific chapter really did get resolved, cleanly enough for the next architectural idea to stand on solid ground, exactly as this book has shown, with working code, at every step. What never got resolved, and what this final chapter's open problems make plain is still live today, data, interpretability, cost, reliability, autonomy, is the question underneath all of them: how much should a system be trusted to generate on its own, and what has to be checked against something outside it before that generation counts as knowledge rather than merely fluent output. The specific architectures in this history all found their own partial answers to specific versions of that question. The question itself, judging by how consistently it has resurfaced across two hundred and fifty years and sixteen chapters, is not one this field, or any other, should expect to finish answering soon.
